\documentclass[pdflatex,sn-mathphys-num]{sn-jnl}

\usepackage{amsmath,amssymb,amsfonts}
\usepackage{mathtools}
\usepackage{booktabs}

\hypersetup{
  pdftitle={Linear Nonequilibrium Information and Conditional Fault Tolerance for Subquantum Search},
  pdfauthor={Michael A. Stay},
  pdfsubject={Foundations of Physics submission},
  pdfkeywords={de Broglie-Bohm theory, quantum nonequilibrium, subquantum computation, quantum error correction, fault tolerance, computational complexity}
}

\raggedbottom
\setcounter{secnumdepth}{3}

\theoremstyle{thmstyleone}
\newtheorem{theorem}{Theorem}
\newtheorem{lemma}{Lemma}
\newtheorem{proposition}{Proposition}
\newtheorem{corollary}{Corollary}

\begin{document}

\title[Conditional Fault Tolerance for Subquantum Search]{Linear Nonequilibrium Information and Conditional Fault Tolerance for Subquantum Search}

\author*[1]{\fnm{Michael A.} \sur{Stay}}\email{mike@cs.auckland.ac.nz}

\affil*[1]{\orgdiv{Computer Science Department}, \orgname{University of Auckland},
\orgaddress{\city{Auckland}, \country{New Zealand}}}

\abstract{
Valentini has proposed that nonequilibrium de Broglie--Bohm systems could enable ``subquantum'' computation, while leaving open how the required nonequilibrium resources scale with input size. For an $n$-bit zero-versus-one-witness problem, let $p_*=2^{-n}$. The initial uniform superposition already contains every candidate basis component; a reversible verifier marks the witness component, if one exists, so the actual marked-state weight is either $0$ or $p_*$. We show that a one-dimensional nonequilibrium pointer prepared using the design value $p_*$ selects the marked outcome with unit probability on a unique-witness instance and never selects it on an ideal no-witness instance. The required relative entropy is $R=\ln(4/p_*)=n\ln2+O(1)$, and a data-processing bound gives a matching $\Omega(n)$ lower bound for constant-probability selection. We then prove a general robustness-transfer lemma: if a conventional implementation has final Born-rule failure probability $\varepsilon$, the same final event has nonequilibrium probability at most $4\varepsilon/p_*$. Thus any fault-tolerant architecture with $\varepsilon_k(n)\le \operatorname{poly}(n)\theta^{2^k}$ requires only $k=O(\log n)$ concatenation levels. This conclusion does not require a classical decomposition into Bohmian fault histories. A five-qubit error-transparent readout illustrates how correctable syndrome sectors can preserve the ideal pointer translation, but a bounded-body scalable implementation is left open. Pointer disturbances are treated separately: an exponential-tail realization gives a constant physical basin margin, so bounded analog errors need not shrink exponentially. Most exploratory derivations and manuscript drafting were carried out interactively with OpenAI ChatGPT (GPT-5.6 Sol), under the author's direction and critical review.
}

\keywords{de Broglie--Bohm theory, quantum nonequilibrium, subquantum computation, quantum error correction, fault tolerance, computational complexity}

\maketitle

\section{Introduction}

In de Broglie--Bohm pilot-wave theory the wave function obeys the Schr\"odinger equation while the configuration follows a deterministic guidance law. Quantum equilibrium corresponds to an ensemble density $\rho(q,t)=|\psi(q,t)|^2$, but the dynamics also permits nonequilibrium ensembles with $\rho\neq |\psi|^2$ \cite{Valentini1991,Valentini2025}. Valentini has emphasized that such ensembles would have operational capabilities absent in ordinary quantum mechanics, including subquantum measurements and discrimination of nonorthogonal states, and proposed that these resources could be used computationally \cite{Valentini2002}.

A recent review makes the missing complexity question explicit: to assess the power of subquantum computation, one must quantify whether the required nonequilibrium resources grow polynomially or exponentially with the size of the computational task \cite{Valentini2025}. We answer a precise version of that question.

The setting is a promise problem with zero or one accepting witness. Write $p_*=2^{-n}$ for the design weight associated with one basis state in the uniform superposition. The superposition already contains every candidate component. A reversible verifier marks the unique witness, when one exists, by correlating its component with a logical flag. Thus the actual marked-state weight is $p_f\in\{0,p_*\}$, rather than a branch being created by the verifier. We prepare a nonequilibrium pointer entirely inside the Bohmian basin leading to the marked outcome when $p_f=p_*$. With relative entropy from quantum equilibrium as the additive resource measure, the resulting cost is linear in $n$ and is optimal up to an additive constant.

The second part treats robustness in two layers. First, we isolate a purely pilot-wave transfer principle. If the initial nonequilibrium density is bounded by a constant multiple $C$ of equilibrium, then common Bohmian transport preserves the event-level inequality $\rho_t(B)\le C\mu_t(B)$ for every later measurable configuration-space event $B$. For the pointer ensemble used here $C=4/p_*$. Consequently a conventional implementation with final Born-rule failure probability $\varepsilon$ has nonequilibrium failure probability at most $4\varepsilon/p_*$. This statement requires no enumeration of Bohmian fault histories. Standard concatenated-code suppression is then used only as one sufficient way of making the conventional failure probability exponentially small. Second, we examine the readout interface itself. A five-qubit error-transparent construction gives exactly the ideal pointer translation throughout the included syndrome sectors, but we do not claim that its concatenated, bounded-body implementation follows automatically from ordinary threshold theorems. That hardware problem is left open. Analog pointer disturbances are handled separately through a constant physical margin in an exponential-tail realization.

Valiant and Vazirani showed that SAT is as hard, under randomized polynomial-time reductions, as distinguishing formulas with no satisfying assignment from formulas with exactly one \cite{ValiantVazirani1986}. A bounded-error polynomial-time solver for the promise problem considered here would therefore yield, after amplification and repeated isolation, a randomized polynomial-time decision procedure for SAT using the same subquantum primitive. A witness could then be recovered by the standard self-reduction of SAT. The complexity consequence is thus randomized, not a deterministic collapse.

\section{Nonequilibrium as an additive information resource}

Let
\begin{equation}
\mu(q)=|\psi(q)|^2
\end{equation}
be the equilibrium measure and let $\rho(q)$ be an actual nonequilibrium ensemble density with $\rho\ll\mu$. Define
\begin{equation}
R(\rho\|\mu)
=\int dq\,\rho(q)\ln\frac{\rho(q)}{\mu(q)}.
\label{eq:resource}
\end{equation}
This is the relative entropy of $\rho$ with respect to $\mu$. It is nonnegative and additive for independent systems. In pilot-wave dynamics, $\rho$ and $|\psi|^2$ obey continuity equations generated by the same velocity field, so the ratio $\rho/|\psi|^2$ is conserved along trajectories \cite{Valentini1991,Valentini2025}. Consequently the fine-grained quantity in Eq.~\eqref{eq:resource} is invariant under exact reversible pilot-wave evolution.

The same quantity gives a lower bound on selection of a small equilibrium set.

\begin{lemma}[Rare-set lower bound]
Let $B$ be a measurable configuration-space event with equilibrium probability $\mu(B)=p$, and let $r=\rho(B)$. If
\begin{equation}
0<p<c\le 1,
\qquad r\ge c,
\end{equation}
then
\begin{equation}
R(\rho\|\mu)\ge
c\ln\frac{c}{p}+(1-c)\ln\frac{1-c}{1-p},
\label{eq:binarylower}
\end{equation}
where the second term is understood as zero when $c=1$. For fixed $c>0$ and $p\rightarrow0$, the right-hand side is $\Omega(\ln(1/p))$.
\end{lemma}

\begin{proof}
Coarse-grain configuration space to the binary random variable specifying membership in $B$. The data-processing inequality gives
\begin{equation}
R(\rho\|\mu)\ge d(r\|p),
\end{equation}
where
\begin{equation}
d(r\|p)=r\ln\frac{r}{p}+(1-r)\ln\frac{1-r}{1-p}.
\end{equation}
For $r\in(p,1)$,
\begin{equation}
\frac{\partial}{\partial r}d(r\|p)
=\ln\frac{r(1-p)}{p(1-r)}>0.
\end{equation}
Because $r\ge c>p$, monotonicity (with the endpoint obtained by continuity when $c=1$) gives $d(r\|p)\ge d(c\|p)$, which is Eq.~\eqref{eq:binarylower}. For fixed $c$,
\begin{equation}
d(c\|p)=c\ln\frac1p+O(1)
\end{equation}
as $p\rightarrow0$.
\end{proof}

In particular, assigning fixed positive probability to an outcome basin whose equilibrium measure is $2^{-n}$ requires $\Omega(n)$ additive nonequilibrium information in this resource measure.

\section{Unique-witness computation as marked-state selection}

Let $f:\{0,1\}^n\to\{0,1\}$ be computable by a polynomial-size reversible verifier, with the promise that either no input satisfies $f$ or exactly one input does. Let $w\in\{0,1\}$ be the number of witnesses and define
\begin{equation}
p_*=2^{-n},
\qquad
p_f=w p_*.
\label{eq:designactual}
\end{equation}
Prepare
\begin{equation}
|\Psi_f\rangle
=\frac{1}{\sqrt{2^n}}\sum_{x\in\{0,1\}^n}|x\rangle|f(x)\rangle.
\label{eq:verifierstate}
\end{equation}
Let $Q=I\otimes|1\rangle\langle1|$ project onto the marked subspace. Then
\begin{equation}
\langle\Psi_f|Q|\Psi_f\rangle=p_f
=\begin{cases}
0,&\text{no witness},\\[3pt]
p_*,&\text{one witness}.
\end{cases}
\label{eq:pweight}
\end{equation}

Couple the flag to a one-dimensional pointer coordinate $y$ with initial packet $g(y)$ using a von Neumann interaction of the form proposed for subquantum measurement by Valentini \cite{Valentini2002,Valentini2025},
\begin{equation}
H_M(t)=a(t)Qp_y.
\label{eq:measurementH}
\end{equation}
Let
\begin{equation}
G(t)=\int_0^t a(s)\,ds.
\end{equation}
When $p_f>0$, write the unmarked component and the marked state as orthogonal normalized states $|\Phi_0\rangle$ and $|\Phi_1\rangle$. When $p_f=0$, the second term below is simply absent. In either case the joint state can be written as
\begin{equation}
|\Psi(t)\rangle
=\sqrt{1-p_f}\,|\Phi_0\rangle g(y)
+e^{i\varphi}\sqrt{p_f}\,|\Phi_1\rangle g(y-G(t)).
\label{eq:pointerstate}
\end{equation}
Let $\mu_y(y)=|g(y)|^2$. Orthogonality gives the pointer density
\begin{equation}
P(y,t)=(1-p_f)\mu_y(y)+p_f\mu_y(y-G(t))
\end{equation}
and current
\begin{equation}
J(y,t)=\dot G(t)p_f\mu_y(y-G(t)).
\end{equation}
The Bohmian pointer velocity is $\dot y=J/P$ wherever $P>0$.

Define the cumulative equilibrium probability
\begin{equation}
C(y,t)=\int_{-\infty}^{y}P(z,t)\,dz.
\end{equation}
Assuming the current vanishes at $-\infty$, the continuity equation gives $\partial_tC=-J$. Therefore a Bohmian trajectory $Y(t)$ satisfies
\begin{equation}
\frac{d}{dt}C(Y(t),t)
=\partial_tC+\dot Y\,\partial_yC
=-J+\frac{J}{P}P=0.
\label{eq:cumulativeconservation}
\end{equation}
Thus the cumulative Born probability associated with the actual pointer position,
\begin{equation}
u=C(Y(t),t),
\end{equation}
is exactly conserved.

Choose the total displacement so that the two final pointer packets are disjoint and the marked-state packet lies to the right. On a unique-witness instance its total Born weight is $p_*$, so its initial Bohmian basin is precisely
\begin{equation}
u>1-p_*.
\label{eq:basin}
\end{equation}
On a no-witness instance $p_f=0$, the current vanishes and no marked-state packet exists. The same pointer preparation can therefore distinguish the two promised cases.

\section{Optimal linear nonequilibrium cost}

Use the design value $p_*=2^{-n}$ independently of which promised case is supplied. Let $u=C(y,0)$ denote the initial cumulative Born probability and define
\begin{equation}
A_*=
\left\{y:1-\frac{p_*}{2}<u<1-\frac{p_*}{4}\right\}.
\label{eq:shell}
\end{equation}
Its equilibrium probability is
\begin{equation}
\mu_y(A_*)=\frac{p_*}{4}.
\end{equation}
Every point of $A_*$ satisfies $u>1-p_*$. Hence every trajectory starting in $A_*$ enters the marked-state packet on a unique-witness instance. On a no-witness instance the packet is absent and the pointer remains in the unshifted outcome.

Condition equilibrium on $A_*$:
\begin{equation}
\rho_{A_*}(y)=\frac{4}{p_*}\mu_y(y)\mathbf 1_{A_*}(y).
\label{eq:conditioned}
\end{equation}
The resource cost is
\begin{equation}
R(\rho_{A_*}\|\mu_y)
=\ln\frac{4}{p_*}
=n\ln2+\ln4.
\label{eq:linearresource}
\end{equation}
Lemma 1, applied to the marked-state basin with any fixed target success probability $c>0$ and sufficiently large $n$ so that $p_*<c$, gives the matching lower bound. Thus the additive nonequilibrium cost of constant-probability selection of a marked state with weight $2^{-n}$ is
\begin{equation}
R=\Theta(n).
\end{equation}

The exponentially small equilibrium probability need not imply exponentially fine spatial localization. Write the initial survival probability as $S(y)=1-C(y,0)$. Over the relevant tail, take
\begin{equation}
S(y)=K e^{-y/\lambda}
\label{eq:exptail}
\end{equation}
for a fixed length scale $\lambda$, with smooth compact-support cutoffs placed outside the interval used below. The marked-state basin $u>1-p_*$ is equivalently $S(y)<p_*$. The near edge of the preparation shell in Eq.~\eqref{eq:shell} has $S=p_*/2$. Therefore its physical distance from the basin boundary $S=p_*$ is
\begin{equation}
m=\lambda\ln2,
\label{eq:pointermargin}
\end{equation}
a constant independent of $n$. The shell itself, between survival probabilities $p_*/2$ and $p_*/4$, also has width $\lambda\ln2$. Its location moves only $O(n)$ into the tail because $-\ln p_*=n\ln2$.

This gives a simple direct robustness statement for analog pointer errors. If the combined effect of preparation and control errors moves the effective pointer coordinate, or the relevant basin boundary, by less than $m$ in the adverse direction, then the outcome cannot change. Such bounded disturbances therefore do not have to shrink exponentially with $n$. If larger pointer disturbances occur with directly controlled nonequilibrium probability at most $\eta$, their contribution can be added as $\eta$ to the logical failure probability below. This $\eta$ is a direct bound on the actual pointer disturbance, not a small Born probability to which the nonequilibrium likelihood-ratio factor is subsequently applied. A displacement sufficient to separate the two final packets requires only $O(n)$ range for the family just described, so bounded interaction strength gives polynomial control time.

\section{An error-transparent encoded pointer interface}

The construction above is vulnerable if ordinary physical errors generate a spurious marked-state population. We therefore couple the pointer to an encoded logical flag through an error-transparent observable and state explicitly where that observable acts.

For concreteness, take the nondegenerate five-qubit perfect code $[[5,1,3]]$, which corrects an arbitrary error on one physical qubit \cite{Laflamme1996}. The same construction works for any fixed nondegenerate distance-three subspace code. Let $\mathcal C$ be its two-dimensional code space and let $P_{\mathcal C}$ be the code projector. Let $E_0=I$ and let $E_1,\ldots,E_{15}$ be the fifteen single-qubit Pauli errors. Nondegeneracy and distance three give
\begin{equation}
P_{\mathcal C}E_a^\dagger E_bP_{\mathcal C}
=\delta_{ab}P_{\mathcal C}.
\label{eq:KLorthogonal}
\end{equation}
Define the partial isometries
\begin{equation}
V_a=E_aP_{\mathcal C}.
\end{equation}
Their ranges $\mathcal C_a=V_a\mathcal C$ are mutually orthogonal syndrome sectors. Define
\begin{equation}
\mathcal K=\bigoplus_{a=0}^{15}\mathcal C_a,
\qquad
P_{\mathcal K}=\sum_{a=0}^{15}V_aV_a^\dagger.
\end{equation}
For the perfect five-qubit code, $16\dim\mathcal C=32$, so $\mathcal K$ is the full five-qubit Hilbert space. For a nonperfect fixed code the formulas below define the operator on $\mathcal K$, and we take it to vanish on $\mathcal K^\perp$.

Let
\begin{equation}
Q_L=|1_L\rangle\langle1_L|
\end{equation}
be the logical marked-state projector on $\mathcal C$. Define the full-system operator
\begin{equation}
\widetilde Q_L
=\sum_{a=0}^{15} V_aQ_LV_a^\dagger.
\label{eq:Qtilde}
\end{equation}
Equation~\eqref{eq:Qtilde} defines $\widetilde Q_L$ on the full physical system Hilbert space. For the five-qubit code the syndrome sectors already exhaust that space; in the general fixed-code version it acts as zero on $\mathcal K^\perp$. From Eq.~\eqref{eq:KLorthogonal},
\begin{equation}
\widetilde Q_L^\dagger=\widetilde Q_L,
\qquad
\widetilde Q_L^2=\widetilde Q_L,
\end{equation}
so it is a bounded self-adjoint projector. It also satisfies the intertwining relation
\begin{equation}
\widetilde Q_LV_a=V_aQ_L.
\label{eq:intertwine}
\end{equation}
The pointer Hamiltonian
\begin{equation}
H_{\mathrm{ET}}(t)=a(t)\widetilde Q_Lp_y
\label{eq:ETH}
\end{equation}
is therefore self-adjoint on the natural domain of $p_y$. This is the measurement analogue of error-transparent evolution: the logical operation is replicated in every included syndrome sector \cite{VyWangJacobs2013}.

Let
\begin{equation}
U(t_2,t_1)=
\exp\!\left[-i\bigl(G(t_2)-G(t_1)\bigr)\widetilde Q_Lp_y\right].
\end{equation}
If one Pauli-basis fault $E_a$ occurs at time $t$ while the state was previously in the code sector, Eq.~\eqref{eq:intertwine} gives
\begin{equation}
U(T,t)E_aU(t,0)P_{\mathcal C}
=E_a e^{-iG(T)Q_Lp_y}P_{\mathcal C}.
\label{eq:transparentidentity}
\end{equation}
Thus the final pointer translation is exactly the ideal logical translation, independent of the identity or timing of this fault. Every operator acting on one physical qubit is a linear combination of the Pauli basis, so by linearity the same identity holds for an arbitrary one-qubit fault until syndrome resolution.

There is also a useful sequential statement. Define an isometry
\begin{equation}
W:\mathcal K\longrightarrow \mathcal C\otimes\mathcal H_{\mathrm{syn}},
\qquad
WV_a|\psi_L\rangle=|\psi_L\rangle\otimes|a\rangle.
\end{equation}
Then
\begin{equation}
W\widetilde Q_LW^\dagger=Q_L\otimes I_{\mathrm{syn}}.
\end{equation}
Consequently, any fault operator $S$ that preserves $\mathcal K$ and obeys
\begin{equation}
WSW^\dagger=I_L\otimes s
\label{eq:syndromeonly}
\end{equation}
commutes with the encoded pointer evolution on $\mathcal K$. Any sequence of such faults at arbitrary times can be moved through the intervening pointer evolutions, so the total pointer translation remains ideal. We call such a history \emph{transparent}. The base gadget includes the no-fault history and every one-qubit fault history covered by Eq.~\eqref{eq:transparentidentity}. We do not assume that an arbitrary multi-fault history remains in $\mathcal K$ or satisfies Eq.~\eqref{eq:syndromeonly}; any history that does not is classified as nontransparent.

Syndrome extraction records which included sector is occupied without learning the logical value. This is the information-theoretic content of the Knill--Laflamme condition: within the correctable error algebra, the syndrome factor is independent of the encoded state \cite{KnillLaflamme1997}. Measurements of stabilizer syndromes also decohere coherent local error components toward stochastic Pauli-like logical noise \cite{Beale2018}.

Equation~\eqref{eq:ETH} is an exact fixed-size base interaction, not by itself a scalable hardware construction. Error-transparent evolution generally requires higher-body interactions \cite{VyWangJacobs2013}. At higher concatenation levels a logical readout acts on an encoded block, and decomposing that readout into bounded-body microscopic operations introduces additional locations whose exact Bohmian transparency is not established by Eq.~\eqref{eq:transparentidentity}. We therefore do not infer a threshold for this particular interface directly from a conventional concatenated-code threshold theorem. Instead, the next section separates the general pilot-wave robustness statement from a conditional fault-tolerance corollary. The corollary applies to any encoded architecture that supplies a conventional Born-rule failure bound while keeping the marked/unmarked logical value opaque to ordinary equilibrium records until the nonequilibrium pointer acquires it. Constructing such an architecture from strictly bounded-body interactions is an open problem.

\section{Robustness transfer and conditional fault tolerance}

We now separate the part specific to pilot-wave nonequilibrium from the conventional fault-tolerance input. Let $\mu_t$ denote the equilibrium measure on the full configuration space at time $t$, including the computer, any environment included in the model, and the pointer. Let $\rho_t$ be the actual nonequilibrium measure transported by the same Bohmian flow.

\begin{lemma}[Event-bound transport]
\label{lem:eventtransport}
Suppose $\rho_0\ll\mu_0$ and
\begin{equation}
\frac{d\rho_0}{d\mu_0}\le C
\qquad \mu_0\text{-almost everywhere}.
\label{eq:initialRN}
\end{equation}
If $\rho_t$ and $\mu_t$ are pushforwards under the same pilot-wave flow, then for every later measurable configuration-space event $B$,
\begin{equation}
\rho_t(B)\le C\mu_t(B).
\label{eq:eventBound}
\end{equation}
Equivalently, $d\rho_t/d\mu_t\le C$ $\mu_t$-almost everywhere.
\end{lemma}

\begin{proof}
Let $\Phi_t$ denote the Bohmian flow, on the full-measure set on which it is defined. Common transport means
\begin{equation}
\rho_t=(\Phi_t)_\#\rho_0,
\qquad
\mu_t=(\Phi_t)_\#\mu_0.
\end{equation}
For a measurable event $B$,
\begin{align}
\rho_t(B)
&=\rho_0(\Phi_t^{-1}(B))\\
&=\int_{\Phi_t^{-1}(B)}\frac{d\rho_0}{d\mu_0}\,d\mu_0\\
&\le C\mu_0(\Phi_t^{-1}(B))\\
&=C\mu_t(B).
\end{align}
The Radon--Nikodym statement follows because the event inequality holds for every measurable $B$. In the usual differential description this is the same fact as advection of the ratio $\rho/|\Psi|^2$ by the Bohmian velocity field.
\end{proof}

For the pointer preparation in Section 4, condition the full equilibrium ensemble on the pointer event $A_*$:
\begin{equation}
\rho_0(q)=\frac4{p_*}\mathbf 1_{A_*}(y(q))\mu_0(q).
\label{eq:fullconditioned}
\end{equation}
Since $\mu_0(A_*)=p_*/4$, this is normalized and
\begin{equation}
R(\rho_0\|\mu_0)=\ln\frac4{p_*},
\qquad
0\le\frac{d\rho_0}{d\mu_0}\le\frac4{p_*}.
\label{eq:ratioBound}
\end{equation}

\begin{theorem}[Equilibrium-to-nonequilibrium robustness transfer]
\label{thm:transfer}
Let $F$ be a measurable final configuration-space event containing every logical or readout failure of an implementation, and suppose its equilibrium Born probability satisfies
\begin{equation}
\mu_T(F)\le\varepsilon.
\end{equation}
For the nonequilibrium ensemble in Eq.~\eqref{eq:fullconditioned},
\begin{equation}
\rho_T(F)\le\frac{4\varepsilon}{p_*}.
\label{eq:transferbound}
\end{equation}
If, in addition, pointer disturbances not represented in $F$ have directly bounded nonequilibrium probability at most $\eta$, then the total failure probability is at most
\begin{equation}
\Pr_{\rho}(\mathrm{failure})
\le \frac{4\varepsilon}{p_*}+\eta.
\label{eq:transfereta}
\end{equation}
\end{theorem}

\begin{proof}
Apply Lemma~\ref{lem:eventtransport} with $C=4/p_*$ to $F$. The second statement is the union bound between $F$ and the separately controlled pointer-disturbance event.
\end{proof}

The theorem deliberately refers only to a final failure event. It does not require physical fault histories to be disjoint Bohmian alternatives, nor does it require a classical decomposition of an unsyndromed coherent bath. All such questions are delegated to whatever conventional analysis supplies the equilibrium probability $\varepsilon$.

As one important sufficient case, suppose a fault-tolerant encoded implementation has $L(n)=\operatorname{poly}(n)$ logical locations and its conventional equilibrium analysis yields
\begin{equation}
\varepsilon_k(n)\le L(n)q_k,
\qquad
q_{j+1}\le A q_j^2,
\qquad
\theta=Aq_0<1,
\label{eq:FTinput}
\end{equation}
where $q_k$ is an upper bound on the level-$k$ logical-location failure probability. The standard concatenated distance-three analysis gives this double-exponential form for suitable fault models below threshold \cite{AliferisGottesmanPreskill2006}; here it is used only as a conventional Born-rule input.

\begin{corollary}[Conditional fault-tolerance scaling]
\label{cor:FT}
Assume an encoded implementation satisfies Eq.~\eqref{eq:FTinput}, and assume that outside its final failure event the computation is logically correct and no ordinary equilibrium record has acquired the marked/unmarked logical value before the nonequilibrium pointer readout. Then for any fixed $\delta\in(0,1)$, the logical contribution to nonequilibrium failure can be made at most $\delta$ with $k=O(\log n)$ concatenation levels. If the cost of one level multiplies each logical location by a constant factor $M$, the total logical fault-tolerance overhead remains polynomial.
\end{corollary}

\begin{proof}
Iteration of the recursion in Eq.~\eqref{eq:FTinput} gives
\begin{equation}
q_k\le\frac1A\theta^{2^k},
\end{equation}
and therefore
\begin{equation}
\varepsilon_k(n)
\le\frac{L(n)}A\theta^{2^k}.
\label{eq:epsilonk}
\end{equation}
Theorem~\ref{thm:transfer} gives
\begin{equation}
\Pr_{\rho}(F_{n,k})
\le\frac{4L(n)}{A p_*}\theta^{2^k}.
\label{eq:neqbad}
\end{equation}
Since $p_*=2^{-n}$, it suffices to choose
\begin{equation}
2^k\ge
\frac{n\ln2+\ln\!\left(4L(n)/(A\delta)\right)}{-\ln\theta}.
\label{eq:kbound}
\end{equation}
For polynomial $L(n)$ the numerator is $O(n)$, so $k=O(\log n)$. The implementation cost is then
\begin{equation}
L(n)M^k=\operatorname{poly}(n).
\end{equation}
The elementary physical error strength in the conventional fault model remains fixed below its threshold; the nonequilibrium resource cost remains $n\ln2+\ln4$.
\end{proof}

The content of Corollary~\ref{cor:FT} is intentionally conditional. The five-qubit construction of Section 5 proves exact transparency for its stated base interaction and fault class, but it does not prove that a bounded-body concatenated implementation of that readout satisfies Eq.~\eqref{eq:FTinput}. Nor should unencoded continuous-pointer faults simply be inserted into $q_0$: concatenating the qubit code does not by itself suppress such faults. Bounded pointer disturbances are instead covered by the constant physical margin in Eq.~\eqref{eq:pointermargin}; rarer larger disturbances may be represented by the separate term $\eta$ in Eq.~\eqref{eq:transfereta} when their probability is controlled directly in the actual nonequilibrium dynamics.

\section{Scope, limitations, and complexity consequence}

Several qualifications are essential.

First, Eq.~\eqref{eq:resource} is an information-theoretic resource measure. It is natural because it is additive and invariant under reversible pilot-wave transport, and because Lemma 1 supplies a matching lower bound. It does not include the unknown physical cost of producing, isolating, or maintaining nonequilibrium matter. The results are conditional on access to the specified pointer ensemble.

Second, Theorem~\ref{thm:transfer} is more general than the earlier fault-history formulation: it needs only a measurable final failure event and its equilibrium Born probability. It does not assume that microscopic fault paths become disjoint classical alternatives. The genuinely physical input is instead that the full equilibrium and nonequilibrium ensembles are transported by the same Bohmian flow, together with a conventional analysis that supplies a bound on the final Born-rule failure probability.

Third, Corollary~\ref{cor:FT} does not establish a threshold for the specific error-transparent readout in Section 5. The constant $A$, the base error parameter $q_0$, and the existence of a bounded-body gadget satisfying Eq.~\eqref{eq:FTinput} would have to be derived for a concrete architecture. In particular, pointer faults and control errors acting on an unencoded continuous coordinate are not automatically suppressed by concatenating the qubit code. The exponential-tail construction instead gives an $n$-independent physical margin for bounded pointer disturbances; larger disturbances require a separate direct bound such as $\eta$.

Fourth, logical opacity remains an architectural requirement. Ordinary equilibrium syndrome or environment records may identify correctable error information, but they must not acquire the marked/unmarked logical value before the nonequilibrium pointer is the first readout carrying that information. The Knill--Laflamme structure supplies this opacity within the correctable syndrome decomposition, but extending it through a complete local fault-tolerant readout is part of the open implementation problem.

Fifth, the error-transparent base operator has been defined exactly on the included syndrome sectors, but we have not derived a scalable realization from a restricted microscopic Hamiltonian. Known error-transparent constructions require higher-body interactions in general \cite{VyWangJacobs2013}. At higher concatenation levels, decomposing encoded readout operations into two-body or other bounded-body primitives introduces additional locations and can spoil exact transparency. Establishing a complete local architecture with both a conventional threshold and the required logical opacity is therefore an open problem rather than a consequence claimed here.

Finally, the unique-witness promise and the randomized nature of the complexity consequence should be kept distinct. Valiant and Vazirani construct randomized polynomial-time isolation reductions with the following properties: if the original SAT formula is unsatisfiable, every derived formula is unsatisfiable; if it is satisfiable, a derived formula has a unique satisfying assignment with inverse-polynomial probability \cite{ValiantVazirani1986}. The promise solver may behave arbitrarily on derived formulas with more than one solution. Repeating the isolation procedure and amplifying the solver's bounded error gives a randomized polynomial-time decision procedure for SAT using the assumed subquantum primitive; with one-sided error on zero instances it gives the corresponding RP-type procedure. This is a statement about randomized polynomial time in the enlarged computational model, not a claim that ordinary BPP contains NP. Once a SAT decision procedure is available, a satisfying assignment can be recovered by the standard polynomial-time self-reduction that fixes variables one at a time. The present note establishes only the resource scaling of the zero-versus-one primitive used in this implication.

\section{Conclusion}

For a marked state with design weight $p_*=2^{-n}$, exponentially small equilibrium probability does not by itself imply an exponential nonequilibrium resource. Conditioning a pilot-wave pointer on an appropriate cumulative-probability shell gives deterministic marked-state selection on the unique-witness case at relative-entropy cost
\begin{equation}
R=n\ln2+O(1),
\end{equation}
and the data-processing argument shows that linear scaling is optimal for constant-probability selection in this resource measure.

The robustness result can be stated without enumerating pilot-wave fault histories. Common Bohmian transport preserves an initial likelihood-ratio bound at the level of every later measurable event. For the pointer ensemble used here, any final event with equilibrium probability $\varepsilon$ therefore has nonequilibrium probability at most $4\varepsilon/p_*$. A conventional fault-tolerant implementation whose Born-rule failure probability decreases as $\operatorname{poly}(n)\theta^{2^k}$ consequently needs only $k=O(\log n)$ concatenation levels to overcome the $O(2^n)$ nonequilibrium amplification. This is a transfer theorem and a conditional corollary, not a threshold proof for the specific readout hardware.

The five-qubit error-transparent construction shows how exact pointer translation can coexist with correctable syndrome information at a fixed base interface. The exponential-tail pointer construction separately shows that bounded analog disturbances can be tolerated with an $n$-independent physical margin. What remains open is the physical synthesis of these ingredients: a strictly bounded-body, scalable readout architecture that simultaneously has a conventional fault-tolerance threshold, keeps the logical marked/unmarked value opaque to ordinary equilibrium records until the nonequilibrium pointer interaction, and controls larger pointer disturbances directly. Thus the information-resource scaling is polynomial under explicit conditions, while the principal unresolved questions are architectural and physical rather than information-theoretic.

\section{Use of generative AI and author contributions}

Most of the exploratory derivation and most of the first-draft manuscript text were produced through an extended interactive use of OpenAI ChatGPT (GPT-5.6 Sol) in September--October 2026. The model proposed and developed several intermediate constructions, including the cumulative-probability-conditioned marked-state pointer argument and the error-transparent fault-tolerance analysis, in response to the author's prompts and critiques. The human author supplied the research question, introduced the branchwise nonlinear-gate idea that motivated the search for a digital marked-state amplifier, repeatedly challenged intermediate claims, directed the successive refinements, selected the scope and claims retained here, and is the sole accountable author. ChatGPT is not listed as an author. The author takes responsibility for the accuracy, integrity, and originality of the submitted manuscript and for verification of the mathematical arguments and references.

\section*{Statements and Declarations}

\noindent\textbf{Funding.} No funding was received to assist with the preparation of this manuscript.

\medskip
\noindent\textbf{Competing interests.} The author declares no competing interests.

\medskip
\noindent\textbf{Data availability.} No datasets were generated or analysed during the current study.

\medskip
\noindent\textbf{Ethics approval and consent to participate.} Not applicable; this work is theoretical and involved no human participants or animals.

\begin{thebibliography}{99}

\bibitem{Valentini2025}
Valentini, A.: De Broglie--Bohm quantum mechanics. In: Encyclopedia of Mathematical Physics, 2nd edn., pp. 24--41 (2025). \url{https://doi.org/10.1016/B978-0-323-95703-8.00264-0}

\bibitem{Valentini1991}
Valentini, A.: Signal-locality, uncertainty, and the subquantum H-theorem. I. Phys. Lett. A \textbf{156}, 5--11 (1991). \url{https://doi.org/10.1016/0375-9601(91)90116-P}

\bibitem{Valentini2002}
Valentini, A.: Subquantum information and computation. Pramana \textbf{59}, 269--277 (2002). \url{https://doi.org/10.1007/s12043-002-0117-1}

\bibitem{ValiantVazirani1986}
Valiant, L.G., Vazirani, V.V.: NP is as easy as detecting unique solutions. Theor. Comput. Sci. \textbf{47}, 85--93 (1986). \url{https://doi.org/10.1016/0304-3975(86)90135-0}

\bibitem{Laflamme1996}
Laflamme, R., Miquel, C., Paz, J.P., Zurek, W.H.: Perfect quantum error correcting code. Phys. Rev. Lett. \textbf{77}, 198--201 (1996). \url{https://doi.org/10.1103/PhysRevLett.77.198}

\bibitem{KnillLaflamme1997}
Knill, E., Laflamme, R.: Theory of quantum error-correcting codes. Phys. Rev. A \textbf{55}, 900--911 (1997). \url{https://doi.org/10.1103/PhysRevA.55.900}

\bibitem{VyWangJacobs2013}
Vy, O., Wang, X., Jacobs, K.: Error-transparent evolution: the ability of multi-body interactions to bypass decoherence. New J. Phys. \textbf{15}, 053002 (2013). \url{https://doi.org/10.1088/1367-2630/15/5/053002}

\bibitem{AliferisGottesmanPreskill2006}
Aliferis, P., Gottesman, D., Preskill, J.: Quantum accuracy threshold for concatenated distance-3 codes. Quantum Inf. Comput. \textbf{6}, 97--165 (2006). \url{https://doi.org/10.26421/QIC6.2-1}

\bibitem{Beale2018}
Beale, S.J., Wallman, J.J., Guti\'errez, M., Brown, K.R., Laflamme, R.: Quantum error correction decoheres noise. Phys. Rev. Lett. \textbf{121}, 190501 (2018). \url{https://doi.org/10.1103/PhysRevLett.121.190501}

\end{thebibliography}

\end{document}
